% Article VII. Authorized additive incorporation, 10-09-2026.
% RESULTADO_RECUPERADO: real Clifford representation and CAR:
% TEORIA_HOLOGRAFICA_INTEGRAL_MASA_HMT_MD_REV2_2026-08-20/fuente/modulos/
% 18_cuantica_estadistica_y_termodinamica.md, lines 131--181 and 218--255;
% first-order action: 14_topologia_exoticos_y_gravedad.md, 459--517;
% distinction between mean current and second moment, and dilution:
% HMT_OBRAS/OBRA_II/chapters/11_cartan_bianchi_cosmologia.tex, 201--230.
% PDF REV2: b56bb38d10b9a4f4d0de47b7da8010fb31a6e8e8684a3bfd4237e37d7a3dddfa.
% FORMALIZACION_NUEVA: explicit affine spinorial action, contraction
% normalized through 31, quartic CAR observable, and amplitude functional.
% CERTIFICADO_NUEVO: rational calculation of the quadratic form in the four
% components of a trivector and evaluation of the finite CAR realization.
% Explicit convention: matrices g have signature (-+++); A_D=-i g^0
% and the symmetric kinetic action has hbar/2. The equivalent combination
% i hbar Gamma^I uses Gamma^I=-i g^I and opposite signature. They are not mixed.
% Dependencies: co-tetrad/connection of 30d; inherited internal constants;
% current and inverses of 31. Material state, cell, and normal ordering are
% data of the stated realization, not extracted universal constants.
% The affine action of this fragment and the minimal-extension action of 30e
% retain their own variational problems. Their currents are not identified.
% 50 is unchanged: its positive-amplitude domain is made concrete.

\section{Spinorial current and torsional amplitude of the material state}
\label{vii:sec:corriente-espinorial-densidad}

The preceding development and the spinorial realization below articulate
two variational problems with common geometric antecedents. The
minimal-extension action \(S_{\min}=\mathfrak a_m E\), constructed in
Section~\ref{vii:sec:corriente-extension-minima}, incorporates the
dependence of transported incidence and its minimum on the
connection; its torsional elimination preserves that complete dependence in
Section~\ref{vii:sec:eliminacion-torsional-no-lineal}. The spinorial
action \(S_D\) in \eqref{vii:eq:accion-dirac-material} couples the
same connection to the fields of its representation and is affine in
contorsion when the co-tetrad and fields are fixed. Each current is obtained
by varying its own action. A minimizing preparation may determine
the initial material state of the spinorial realization; substituting
connection-dependent fields also retains the terms of the
composed differential. This distinction places the following amplitude
calculation within its specific action and variational domain.

The spinorial representation of transport and the algebraic
Cartan--Holst response allow torsional intensity to be calculated as a
functional of the material state. Its evaluation preserves three data:
the field representation, the second moment of its current, and
the volume with respect to which density is defined. The coupling
coefficient follows from the action and its inversion; the initial state
determines the material amplitude.

\subsection{Spinorial representation and affine material action}
\label{vii:subsec:accion-espinorial-afin}

We retain the co-tetrad, orientation, and couplings of
Section~\ref{vii:sec:corriente-respuesta-cartan}. The spinorial
realization uses the four-dimensional complex module obtained by
complexifying the real matrices
\begin{equation}
 \begin{gathered}
 J=\begin{pmatrix}0&-1\\1&0\end{pmatrix},\qquad
 R=\begin{pmatrix}0&1\\1&0\end{pmatrix},\qquad
 Z=\begin{pmatrix}1&0\\0&-1\end{pmatrix},\\
 g^0=J\otimes I_2,\qquad g^1=R\otimes I_2,\qquad
 g^2=Z\otimes R,\qquad g^3=Z\otimes Z.
 \end{gathered}
 \label{vii:eq:clifford-material-explicito}
\end{equation}
Multiplication gives
\(\{g^I,g^J\}=2\eta^{IJ}I_4\), with
\(\eta=\operatorname{diag}(-1,1,1,1)\). The matrices are distinguished
from the scalar parameter \(\gamma\) of the Holst operator. Define
\begin{equation}
 g_I=\eta_{IJ}g^J,\quad
 \mathbb B_{IJ}=\tfrac14[g_I,g_J],\quad
 A_D=-i g^0,\quad \bar\psi=\psi^\dagger A_D,\quad
 \Omega=\tfrac12\omega^{IJ}\mathbb B_{IJ}.
 \label{vii:eq:adjunto-conexion-espinorial}
\end{equation}
The Clifford product satisfies
\[
 [\mathbb B_{IJ},g^K]=\delta_J^K g_I-\delta_I^K g_J.
\]
Thus the same metric connection acts on the fields through
\(D\psi=\mathrm d\psi+\Omega\psi\) and
\(D\bar\psi=\mathrm d\bar\psi-\bar\psi\Omega\).
Transport along a route is obtained by integrating this connection in
its spinorial representation; it preserves the connection and route of
the preceding geometric realization.

Let \(\eta_I=\iota_{E_I}\operatorname{vol}_e\), so that
\(e^J\wedge\eta_I=\delta_I^J\operatorname{vol}_e\). For a
field of dimensional type \([\psi]=L^{-3/2}\), take the minimal
material action
\begin{equation}
 S_D=\int_M\left\{
 \frac{\hbar}{2}
 \bigl[\bar\psi g^I D\psi-(D\bar\psi)g^I\psi\bigr]\wedge\eta_I
 -mc\,\bar\psi\psi\operatorname{vol}_e\right\}.
 \label{vii:eq:accion-dirac-material}
\end{equation}
Here \(\hbar=\hbar_{\rm int}\), \(c=c_{\rm int}\), and the mass
\(m\) is the reading of the material species in the chosen chart.
These values are outputs preceding this realization. The matrix
\(A_D\) is Hermitian and \(A_Dg^I\) is anti-Hermitian. The
symmetric kinetic term therefore has the real convention corresponding to
\eqref{vii:eq:clifford-material-explicito}. In this signature, the
torsion-free equation has principal term \(\hbar g^I D_I\); the equivalent
notation \(i\hbar\Gamma^I D_I\), with \(\Gamma^I=-i g^I\),
uses the Clifford relation of opposite signature. This precision fixes
the adjoint, signature, and imaginary factor of the action jointly.

The material fields and co-tetrad remain fixed during connection
variation. A preparation from the minimal extension in
Section~\ref{vii:sec:corriente-extension-minima} may determine
the initial field or state datum. Each action retains its own
variation: if a connection-dependent minimizing field is substituted
into an action, its composed differential also applies,
or the envelope theorem under its stationarity hypotheses.
Action~\eqref{vii:eq:accion-dirac-material} realizes here a
coupling affine in contorsion.

\begin{proposition}[Spinorial current with variational normalization]
\label{vii:prop:corriente-dirac-explicita}
Let
\(B^{IJK}=\bar\psi g^{[I}g^Jg^{K]}\psi\), where brackets
denote normalized antisymmetrization. The current defined in
\eqref{vii:eq:corriente-variacional} by action
\eqref{vii:eq:accion-dirac-material} is
\begin{equation}
 \sigma_{JK}
 =-\frac{\hbar}{2}\,
   \bar\psi\{g^I,\mathbb B_{JK}\}\psi\,\eta_I
 =-\frac{\hbar}{2}\,B^I{}_{JK}\eta_I.
 \label{vii:eq:corriente-dirac-trivector}
\end{equation}
It is independent of contorsion when \(e\) and \(\psi\) are
fixed. The inversion in Section~\ref{vii:subsec:reconstruccion-torsion}
therefore applies to this current.
\end{proposition}
\begin{proof}
The variation of \(\Omega\) is
\(\delta\Omega=\frac12\delta\omega^{JK}\mathbb B_{JK}\).
The two kinetic summands give, respectively, the product with
\(g^I\) on the left and on the right. Consequently,
\[
 \delta_\omega S_D
 =\frac{\hbar}{4}\int_M\delta\omega^{JK}\wedge\eta_I\,
                 \bar\psi\{g^I,\mathbb B_{JK}\}\psi.
\]
Comparison with \eqref{vii:eq:corriente-variacional} fixes the
factor and sign in \eqref{vii:eq:corriente-dirac-trivector}.
Expanding the anticommutator and using the Clifford relation
cancels the degree-one contractions, leaving
\(\{g^I,\mathbb B_{JK}\}=g^{[I}g_Jg_{K]}\).
The action is linear in \(\Omega\), which proves the final
statement.
\end{proof}

\subsection{Effective contraction and second-moment coefficient}
\label{vii:subsec:contraccion-espinorial}

The Lorentzian quadratic form of the trivector is written
\begin{equation}
 q(B)=\frac1{3!}B^{IJK}B_{IJK}
 =-(B^{012})^2-(B^{013})^2-(B^{023})^2+(B^{123})^2.
 \label{vii:eq:cuadratica-trivector}
\end{equation}
This contraction preserves the signature of the spinorial module. Its second
moment depends on the state and on the realization of field
products.

\begin{theorem}[Effective interaction of the spinorial current]
\label{vii:thm:coeficiente-torsional-espinorial}
With material action~\eqref{vii:eq:accion-dirac-material}, algebraic
elimination of contorsion gives
\begin{equation}
 \mathcal L_{\rm spin}^{\rm eff}
 =C_\gamma q(B)\operatorname{vol}_e,\qquad
 C_\gamma=\frac{3\kappa c\hbar^2}{16}
                      \frac{\gamma^2}{1+\gamma^2}.
 \label{vii:eq:coeficiente-cuadratico-espinorial}
\end{equation}
The coefficient is obtained from the current and the Holst inverse
with the normalizations of the preceding action.
\end{theorem}
\begin{proof}
By \eqref{vii:eq:accion-efectiva-spin}, it suffices to evaluate
\(-\frac14K_*^{IJ}\wedge\sigma_{IJ}\), with
\[
 Y^{IJ}=\kappa c\frac{\gamma^2}{1+\gamma^2}
           (-\star+\gamma^{-1}I)\sigma^{IJ},\qquad
 K_*=\mathsf C_e(Y).
\]
Here \(\mathsf C_e\) is the explicit composition of
\eqref{vii:eq:torsion-explicita} with
\eqref{vii:eq:contorsion-componentes}. Linearity of these two
inverses reduces the calculation to the \(-\star\) and \(I\) parts.
To give the complete contraction, order the trivector components
as \((012,013,023,123)\), factor out
\(\kappa c\hbar^2\gamma^2/(1+\gamma^2)\), and substitute
\(\sigma_{JK}=-\frac12B^I{}_{JK}\eta_I\). The two matrices
of the remaining quadratic form are
\begin{equation}
 \mathcal B_{-\star}
 =\frac3{16}\operatorname{diag}(-1,-1,-1,1),\qquad
 \mathcal B_I=0_{4\times4}.
 \label{vii:eq:matrices-contraccion-espinorial}
\end{equation}
These matrices are obtained using only
\(e^I\wedge\eta_J=\delta_J^I\operatorname{vol}_e\): for
each component \(B^{abc}\), form \(\sigma\), compute
\(Y\), then
\(T^I=\sum_J\iota_{E_J}Y^{IJ}
 -\frac14e^I\wedge\sum_{L,J}\iota_{E_L}\iota_{E_J}Y^{LJ}\),
and finally
\(K_{IJK}=(T_{KIJ}-T_{IJK}-T_{JKI})/2\).
Contraction gives the four displayed diagonal entries; the
six polarizations between distinct components are zero for both
operators. This calculates every coefficient of the two
quadratic forms. The \(\gamma^{-1}I\) part vanishes, and the
\(-\star\) part produces exactly \eqref{vii:eq:cuadratica-trivector},
with the coefficient in \eqref{vii:eq:coeficiente-cuadratico-espinorial}.
\end{proof}

\subsection{Finite fermionic realization and the second moment of the state}
\label{vii:subsec:segundo-momento-car}

Fix a periodic comoving cell, an orthonormal frame adapted to its
observer, and the homogeneous spatial mode of the four spinorial
components. This finite realization has state space
\(\mathcal F=\bigoplus_{n=0}^4\Lambda^n\mathbb C^4\).
Let \(a_r,a_r^\dagger\), \(r=1,\ldots,4\), be its canonical
operators, with
\(\{a_r,a_s^\dagger\}=\delta_{rs}\) and
\(\{a_r,a_s\}=0\). The occupation vacuum fixes normal
ordering. For a matrix \(M\), write
\[
 \mathrm d\Gamma(M)=\sum_{r,s}a_r^\dagger M_{rs}a_s,
 \qquad \widehat N=\mathrm d\Gamma(I_4).
\]
The letter \(\Gamma\) in this notation denotes second
quantization; nonadic holonomy retains its own notation.

The four matrices of the bilinears are calculated directly:
\begin{equation}
 \begin{array}{c|cccc}
 abc&012&013&023&123\\ \hline
 M_{abc}=A_Dg^ag^bg^c
       &iJ\otimes R&iJ\otimes Z&iI_2\otimes J&iZ\otimes J
 \end{array}
 \label{vii:eq:matrices-bilineales-espin}
\end{equation}
All are Hermitian, have zero trace, and satisfy
\(M_{abc}^2=I_4\). The quartic observable corresponding to
\eqref{vii:eq:cuadratica-trivector}, with this stated normal ordering,
is
\begin{equation}
 \begin{split}
 \widehat{\mathscr Q}_{\rm sp}
 &=\sum_{a<b<c}\eta_{aa}\eta_{bb}\eta_{cc}
       :\!\mathrm d\Gamma(M_{abc})^2\!:\\
 &=\sum_{a<b<c}\eta_{aa}\eta_{bb}\eta_{cc}
       \mathrm d\Gamma(M_{abc})^2+2\widehat N.
 \end{split}
 \label{vii:eq:observable-cuadratico-espin}
\end{equation}
Normal ordering is part of the specified finite realization;
its definition makes explicit the subtraction of single-occupation
contractions.

\begin{proposition}[Evaluation by the canonical relations]
\label{vii:prop:operador-espin-car}
The operator \(\widehat{\mathscr Q}_{\rm sp}\) is self-adjoint,
preserves each occupation sector, and has restrictions
\begin{equation}
 \widehat{\mathscr Q}_{\rm sp}|_{\Lambda^0\oplus\Lambda^1}=0,
 \qquad
 \widehat{\mathscr Q}_{\rm sp}|_{\Lambda^3}=4I_4,
 \qquad
 \widehat{\mathscr Q}_{\rm sp}|_{\Lambda^4}=8I_1.
 \label{vii:eq:sectores-espin-car}
\end{equation}
In the basis of \(\Lambda^2\mathbb C^4\) whose occupation
subsets are \((12,13,23,14,24,34)\), its matrix is
\begin{equation}
 \begin{pmatrix}
 0&0&0&0&0&-4\\
 0&2&0&0&6&0\\
 0&0&2&-6&0&0\\
 0&0&-6&2&0&0\\
 0&6&0&0&2&0\\
 -4&0&0&0&0&0
 \end{pmatrix}.
 \label{vii:eq:sector-dos-espin-car}
\end{equation}
\end{proposition}
\begin{proof}
The relation \(a_s a_t^\dagger=\delta_{st}-a_t^\dagger a_s\)
applied to \(\mathrm d\Gamma(M)^2\) separates its single-occupation
contraction and gives
\[
 :\!\mathrm d\Gamma(M)^2\!:
 =\mathrm d\Gamma(M)^2-\mathrm d\Gamma(M^2).
\]
Since the first three signs in
\eqref{vii:eq:matrices-bilineales-espin} are negative and the last
positive, the contractions sum to \(-2\widehat N\).
This proves \eqref{vii:eq:observable-cuadratico-espin}.
Self-adjointness and preservation of occupation follow from
the same expressions.

On \(\Lambda^0\), all summands vanish. On \(\Lambda^1\),
\(\mathrm d\Gamma(M)=M\), and the squares sum to \(-2I_4\),
which cancels with \(2\widehat N\). On \(\Lambda^4\),
\(\mathrm d\Gamma(M)=\operatorname{tr}M=0\), leaving eight.
The identification by complements between \(\Lambda^3\mathbb C^4\)
and the dual of \(\mathbb C^4\) takes
\(\mathrm d\Gamma(M)\) to
\((\operatorname{tr}M)I-M^{\mathsf T}=-M^{\mathsf T}\).
Their squares again sum to \(-2I_4\), whereas
\(2\widehat N=6I_4\); the restriction is therefore \(4I_4\).
Finally, applying the four explicit matrices to the
six exterior products of the indicated basis produces
\eqref{vii:eq:sector-dos-espin-car}. The operator is thus evaluated
on all sixteen dimensions of \(\mathcal F\).
\end{proof}

For a state operator \(\varrho\geq0\),
\(\operatorname{Tr}\varrho=1\), define
\(q_\varrho=\operatorname{Tr}(\varrho\widehat{\mathscr Q}_{\rm sp})\).
Its dependence on second moments is explicit:
\begin{equation}
 q_\varrho=\sum_{a<b<c}\eta_{aa}\eta_{bb}\eta_{cc}
 \left\{
  \langle\widehat B_{abc}\rangle_\varrho^2
  +\operatorname{Var}_\varrho(\widehat B_{abc})
  -\langle\widehat N\rangle_\varrho
 \right\},\qquad
 \widehat B_{abc}=\mathrm d\Gamma(M_{abc}).
 \label{vii:eq:covarianza-espin-material}
\end{equation}
The covariance and normal-ordering subtraction are preserved even when
the mean currents vanish. In particular, if \(P_3\) is the
projection onto \(\Lambda^3\mathbb C^4\),
\begin{equation}
 \varrho_3=\tfrac14P_3,\qquad
 \langle\widehat B_{abc}\rangle_{\varrho_3}=0,\qquad
 q_{\varrho_3}=4>0.
 \label{vii:eq:testigo-positivo-espin}
\end{equation}
This is a positive trace-one state and provides an explicit witness
for the positive domain of the functional. The two-occupation sector
also contains negative values; for example, the vector
\((|12\rangle+|34\rangle)/\sqrt2\) has value \(-4\).
The sign of the intensity thus originates in the contraction and
the state, and is kept separate from the selection of cosmological conditions.

\subsection{Density, pressure, and preservation of comoving amplitude}
\label{vii:subsec:amplitud-espinorial-comovil}

In the spatially homogeneous realization, fix
\[
 e^0=N(t)c\,\mathrm dt,\qquad e^j=a(t)\,\mathrm dx^j,
 \qquad \mathcal V(t)=a(t)^3V_c,
\]
with lapse \(N(t)>0\), dimensionless scale factor \(a(t)>0\),
and comoving volume \(V_c>0\). The frame and periodic mode of
Subsection~\ref{vii:subsec:segundo-momento-car} remain specified.
The normalization of the homogeneous field is
\(\widehat\psi=\mathcal V^{-1/2}(a_1,a_2,a_3,a_4)^{\mathsf T}\).
Each bilinear carries a factor \(\mathcal V^{-1}\), and its normally
ordered product carries \(\mathcal V^{-2}\). The symmetric temporal
kinetic form preserves this canonical normalization: the terms arising
from differentiating \(\mathcal V^{-1/2}\) cancel between
its two summands.

\begin{theorem}[Torsional amplitude as a functional of the state]
\label{vii:thm:amplitud-material-explicita}
Holding the canonical state and the product realization fixed
during metric variations, the effective spin term has reduced
action
\begin{equation}
 S_{\rm spin}^{\rm hom}
 =\int\!\mathrm dt\,N(t)
       \frac{cC_\gamma}{a(t)^3V_c}\,q_\varrho.
 \label{vii:eq:accion-homogenea-espin}
\end{equation}
Its energy density and isotropic pressure are
\begin{equation}
 \rho_{\rm spin}=p_{\rm spin}
 =-\frac{cC_\gamma}{V_c^2}\,q_\varrho a^{-6}.
 \label{vii:eq:densidad-presion-espin}
\end{equation}
In any evolution preserving \(q_\varrho\), the corresponding
constant amplitude is
\begin{equation}
 \boxed{\displaystyle
 s_0[\varrho,V_c]
 =\frac{3\kappa c^2\hbar^2}{16}
    \frac{\gamma^2}{1+\gamma^2}
    \frac{\operatorname{Tr}
       (\varrho\widehat{\mathscr Q}_{\rm sp})}{V_c^2}.}
 \label{vii:eq:amplitud-s0-espinorial}
\end{equation}
\end{theorem}
\begin{proof}
Spatial integration of
\eqref{vii:eq:coeficiente-cuadratico-espinorial} multiplies the
factor \(\mathcal V^{-2}\) by
\(Nc\mathcal V\,\mathrm dt\), giving
\eqref{vii:eq:accion-homogenea-espin}. In canonical variables,
lapse variation defines energy through
\(\delta_N S_{\rm m}=-\int E\,\delta N\,\mathrm dt\).
Therefore,
\[
 E_{\rm spin}=-\frac{cC_\gamma q_\varrho}{a^3V_c},
 \qquad \rho_{\rm spin}=\frac{E_{\rm spin}}{a^3V_c}.
\]
The isotropic spatial variation satisfies
\(\delta_a S_{\rm m}=\int3Na^2V_c p\,\delta a\,\mathrm dt\).
Applied to \eqref{vii:eq:accion-homogenea-espin}, it gives
\[
 3Na^2V_c p_{\rm spin}
 =-\frac{3NcC_\gamma q_\varrho}{a^4V_c}.
\]
This proves simultaneously the sign, pressure, and density in
\eqref{vii:eq:densidad-presion-espin}. When
\(q_\varrho\) is preserved, the coefficient of \(-a^{-6}\) is constant,
yielding \eqref{vii:eq:amplitud-s0-espinorial}. The dimensions
are consistent: \([cC_\gamma]=\mathrm{energy}\,L^3\),
so \(s_0\) has the dimension of energy density.
\end{proof}

\Needspace{27\baselineskip}
\begin{corollary}[Preserved domain of positive amplitude]
\label{vii:cor:dominio-positivo-conservado}
Let a unitary evolution of \(\mathcal F\) preserve
\(\Lambda^3\mathbb C^4\). For every trace-one initial state
supported on that sector, one has \(q_{\varrho(t)}=4\), and
\begin{equation}
 s_0=\frac{3\kappa c^2\hbar^2}{4V_c^2}
                   \frac{\gamma^2}{1+\gamma^2}>0.
 \label{vii:eq:amplitud-positiva-sector-tres}
\end{equation}
\end{corollary}
\begin{proof}
Sector invariance preserves the support of the state; on that sector,
\eqref{vii:eq:sectores-espin-car} gives
\(\widehat{\mathscr Q}_{\rm sp}=4I\). Trace preservation
then fixes its expectation at four, independently of the internal
evolution of its coherences. Positivity of the couplings,
\(V_c>0\), and \(\gamma\ne0\) proves the formula
and its sign.
\end{proof}

For general states, preservation of
\(\langle\widehat N\rangle\) determines dilution of the occupation
density, \(n(t)=\langle\widehat N\rangle/(a^3V_c)\).
Preservation of the quartic moment is checked separately.
If \(i\hbar\dot\varrho=[H(t),\varrho]\), then
\begin{equation}
 \frac{\mathrm d q_\varrho}{\mathrm dt}
 =\frac{i}{\hbar}\operatorname{Tr}
       \bigl(\varrho[H(t),\widehat{\mathscr Q}_{\rm sp}]\bigr).
 \label{vii:eq:conservacion-momento-cuartico}
\end{equation}
Commutation, invariance of a sector on which the observable is
scalar, or a conservation law proved for the specific state
provides the relevant conditions. If the moment evolves,
\eqref{vii:eq:densidad-presion-espin} retains its instantaneous
evaluation; its balance includes
\[
 \dot\rho_{\rm spin}+6H\rho_{\rm spin}
 =-\frac{cC_\gamma}{V_c^2}\,a^{-6}\dot q_\varrho,
 \qquad H=\dot a/a
\]
in proper time. This term records exchange with the remaining
material degrees of freedom.

The composition is determined: spinorial representation of
transport, material action, variational current, Cartan--Holst
inversion, second moment, and normalized density. The domain
\(s_0>0\) in Section~\ref{vii:sec:dilucion-rebote} thus has
an evaluable functional and explicit witness states. Its use in
cosmological dynamics also preserves the material, geometric,
and conservation conditions stated in that section.
